\chapter{Introduction}
\label{ch:introduction}

\section{Motivation}
\label{sec:intro-motivation}

Large language models perform well on many natural language tasks, but their
performance is not the same across languages. Models trained mostly on English and
a few other high-resource languages do worse on low-resource languages, even
though the same parameters are used for all of them
\citep{conneau2018xnli, shi2022mgsm}. A second issue is that it is hard to say
which part of a model is responsible for a given behaviour, and equally hard to
change that behaviour by editing a specific part. An intervention that looks
targeted often has little effect on the actual task.

Two lines of research address these issues from different directions. One studies
the model from the inside and asks where a capability is stored. Work on
language-specific neurons follows this direction, suggesting that a small set of
neurons accounts for much of how a model handles a particular language, and that
these neurons can be found from activation statistics
\citep{tang2024language, kojima2024multilingual}. The other line works through the
training objective, using preference-based and reinforcement-learning methods to
push model outputs towards a desired signal
\citep{rafailov2023direct, shao2024deepseekmath}. This thesis looks at both
directions together. The common observation across the studies reported here is
that editing a localized part of a model is often less useful than changing the
objective the model is trained on.

The problem matters for two reasons. In practice, low-resource languages and
imbalanced label sets are settings where data is limited, so a method that
actually moves a model's behaviour, rather than only appearing to, is worth
having. From a research point of view, the cases where an intervention fails are
informative: a neuron edit that changes nothing, or a training signal that moves a
representation without improving the task, tells us something about how these
models are organized.

\section{Problem Statement}
\label{sec:intro-problem}

This thesis studies the gap between what a model can do and what its training
makes it do, in two settings: cross-lingual learning in multilingual language
models, and optimization with competing objectives in multi-label classification.
We break this into three sub-problems.

\begin{enumerate}
	\item \textbf{Can localized internal components be used to control
		cross-lingual behaviour?} If language-specific neurons encode a language in a
	separable way \citep{tang2024language, kojima2024multilingual}, then editing
	them should improve performance on that language. The problem is to test this
	on downstream tasks, not only on language of generation, and to see whether
	the interventions give real gains.
	
	\item \textbf{How does reinforcement-learning fine-tuning change cross-lingual
		behaviour, and why does it help some languages but not others?} Training a
	multilingual model with a group-relative reinforcement objective
	\citep{shao2024deepseekmath} improves some languages and not others. The
	problem is to describe this difference and find the factor that controls it,
	rather than attributing it loosely to how low-resource a language is.
	
	\item \textbf{How can a model be improved on a chosen subset of labels without
		hurting the rest?} In multi-label classification, rare or important labels are
	not well served by a uniform loss, but improving them naively can degrade
	frequent labels. The problem is to design an objective that targets the chosen
	labels while protecting the others.
\end{enumerate}

\section{Research Questions}
\label{sec:intro-questions}

From these sub-problems, the thesis investigates the following questions.

\begin{itemize}
	\item \textbf{Are language-specific neurons useful for cross-lingual
		transfer?} Earlier work shows these neurons can control the \emph{language} a
	model generates in \citep{tang2024language, kojima2024multilingual}. It is not
	clear whether editing them improves performance on tasks such as natural
	language inference \citep{conneau2018xnli} and question answering
	\citep{artetxe2020cross} for low-resource languages.
	
	\item \textbf{Does a multilingual model process its reasoning in English at
		middle layers, and is that what holds back low-resource languages?} One
	hypothesis is that multilingual transformers work in an English-leaning space
	at middle layers \citep{wendler2024llamas}. We ask whether this holds during
	reinforcement-learning fine-tuning, and whether this or some other factor
	explains why certain languages do not improve.
	
	\item \textbf{Can preference optimization, which was developed for text
		generation, be used to manage competing objectives in a classification
		setting?} Methods such as DPO \citep{rafailov2023direct}, CPO
	\citep{xu2024contrastive}, and SimPO \citep{meng2024simpo} were built for
	generation. We ask whether their main ideas, together with a group-robust
	balancing method \citep{ramesh2024group}, can give targeted and fair
	improvements in multi-label classification.
\end{itemize}

\section{Research Details}
\label{sec:intro-details}

The thesis contains three studies, one for each question above. We describe them
briefly here and in full in the later chapters.

\subsection{Language-Specific Neurons and Cross-Lingual Transfer}
\label{sec:intro-neurons}

The first study, published at the Insights from Negative Results workshop at
NAACL~2025 \citep{mondal2025language}, tests whether language-specific neurons can
be used to improve cross-lingual transfer. We identify these neurons in Llama~3.1
and Mistral Nemo using Language Activation Probability Entropy and
activation-probability thresholding \citep{tang2024language}, and we try two kinds
of intervention: editing neuron activations at test time, and fine-tuning only
those neurons with a low-rank adapter \citep{hu2022lora}, on XNLI
\citep{conneau2018xnli} and XQuAD \citep{artetxe2020cross}. The interventions give
gains of less than one accuracy point on low-resource languages. We also find that
setting a neuron's activation to zero does not turn off the behaviour assigned to
it, and that the neurons found for different languages are not independent of one
another. These results do not support a clean, component-level account of how a
model handles language, and are consistent with a distributed encoding in which
single neurons carry more than one function \citep{elhage2022superposition}.

\subsection{Cross-Lingual Behaviour during Policy Optimization}
\label{sec:intro-grpo}

The second study looks at how cross-lingual behaviour changes during
reinforcement-learning fine-tuning, rather than in a fixed model. We fine-tune
Qwen2.5-7B-Instruct on the MGSM benchmark \citep{shi2022mgsm} using Group Relative
Policy Optimization \citep{shao2024deepseekmath} with a low-rank adapter, and we
measure the model at regular checkpoints. The measurements cover task accuracy,
the language of the model's intermediate predictions read off with the logit lens,
and the geometry of the residual-stream representations compared with centred
kernel alignment. Accuracy changes little across the seven languages over
training, and Telugu stays far below the rest throughout. Looking inside the
model, we find that it predicts in English through the early and middle layers and
surfaces the target language only in the last few layers, in line with the
latent-language hypothesis \citep{wendler2024llamas}. This pattern is already
present in the base model and does not change over training, so the fine-tuning
does not cause cross-lingual collapse. Telugu, the language that stays below the
rest, also has the fewest dedicated tokens in the Qwen2.5 tokeniser and is the one
whose representations change the most while its accuracy does not move. We trace
the per-language outcome to tokeniser coverage rather than to language identity:
when the vocabulary cannot represent a language at a fine enough level, the reward
signal is too coarse to give a useful gradient.

\subsection{Robust Preference Optimization for Multi-Label Classification}
\label{sec:intro-fairpo}

The third study, published at the OPT workshop at NeurIPS~2025
\citep{mondal2025fairpo}, uses preference optimization as a tool rather than as an
object of study. We propose FairPO, a method for multi-label classification that
splits the labels into a privileged group and a non-privileged group. The
privileged group is trained with a preference-based objective, adapted from DPO
\citep{rafailov2023direct}, CPO \citep{xu2024contrastive}, and SimPO
\citep{meng2024simpo}, that makes the model prefer a true label over the incorrect
labels most likely to be confused with it. The non-privileged group is kept close
to a reference model by a constrained objective, which prevents its performance
from dropping. The two objectives pull against each other, and we balance them
with a Group Robust Preference Optimization formulation \citep{ramesh2024group}
set up as a minimax problem over the groups. On MS-COCO and NUS-WIDE, FairPO
improves accuracy on rare privileged labels and leaves frequent labels almost
unchanged, doing better than the baselines we compare against.

We note one point about naming. Two different methods in this thesis share the
abbreviation GRPO. In Section~\ref{sec:intro-grpo} and
Chapter~\ref{ch:grpo-crosslingual}, GRPO means Group Relative Policy Optimization
\citep{shao2024deepseekmath}, a reinforcement-learning algorithm. In
Section~\ref{sec:intro-fairpo} and Chapter~\ref{ch:fairpo}, the balancing method
is Group Robust Preference Optimization \citep{ramesh2024group}, which comes from
the robust optimization literature. We write out the full name the first time it
appears in each chapter and do not use the abbreviation on its own to tell them
apart.

\section{Thesis Organization}
\label{sec:intro-organization}

The rest of the thesis is organized as follows.
Chapter~\ref{ch:related-work} reviews related work on multilingual
representations, language-specific neurons, preference optimization, and robust
optimization. Chapter~\ref{ch:neurons} presents the study of language-specific
neurons and cross-lingual transfer. Chapter~\ref{ch:grpo-crosslingual} presents
the study of cross-lingual behaviour during Group Relative Policy Optimization.
Chapter~\ref{ch:fairpo} presents FairPO and its results on multi-label
benchmarks. Chapter~\ref{ch:conclusion} summarizes the findings, states their
limitations, and describes directions for future work.